%-------------------------
% Resume in Latex
% Author : Sourabh Bajaj
% Modified by: Vadim Kholodilo
% License : MIT
%------------------------

\documentclass[letterpaper,11pt]{article}

\usepackage[T2A]{fontenc}
\usepackage[utf8]{inputenc}
\usepackage[russian]{babel}
%-------------------------
% Shared preamble and macros for resume-en.tex and resume-ru.tex
% Each language file sets \documentclass and language packages
% (fontenc/inputenc/babel) before \input-ing this file.
%------------------------

\usepackage{latexsym}
\usepackage[empty]{fullpage}
\usepackage{titlesec}
\usepackage{marvosym}
\usepackage[usenames,dvipsnames]{color}
\usepackage{verbatim}
\usepackage{enumitem}
\usepackage[hidelinks]{hyperref}
\usepackage{fancyhdr}
\usepackage{tabularx}
\input{glyphtounicode}

\pagestyle{fancy}
\fancyhf{} % clear all header and footer fields
\fancyfoot{}
\renewcommand{\headrulewidth}{0pt}
\renewcommand{\footrulewidth}{0pt}

% Adjust margins
\addtolength{\oddsidemargin}{-0.5in}
\addtolength{\evensidemargin}{-0.5in}
\addtolength{\textwidth}{1in}
\addtolength{\topmargin}{-.5in}
\addtolength{\textheight}{1.0in}

\urlstyle{same}

\raggedbottom
\raggedright
\setlength{\tabcolsep}{0in}

% Sections formatting
\titleformat{\section}{
  \vspace{-4pt}\scshape\raggedright\large
}{}{0em}{}[\color{black}\titlerule \vspace{-5pt}]

% Ensure that generate pdf is machine readable/ATS parsable
\pdfgentounicode=1

%-------------------------
% Custom commands
\newcommand{\resumeItem}[1]{
  \item\small{
    \textbf{#1 \vspace{-3pt}}
  }
}

% Per-position tech stack line: \resumeSkills{Label}{comma-separated list}
\newcommand{\resumeSkills}[2]{
  \item[]\small{\textbf{#1: }#2}\vspace{-4pt}
}

% Just in case someone needs a heading that does not need to be in a list
\newcommand{\resumeHeading}[4]{
    \begin{tabular*}{0.99\textwidth}[t]{l@{\extracolsep{\fill}}r}
      \textbf{#1} & #2 \\
      \textit{\small#3} & \textit{\small #4} \\
    \end{tabular*}\vspace{-5pt}
}

\newcommand{\resumeSubheading}[4]{
  \vspace{-1pt}\item
    \begin{tabular*}{0.97\textwidth}[t]{l@{\extracolsep{\fill}}r}
      \textbf{#1} & #2 \\
      \textit{\small#3} & \textit{\small #4} \\
    \end{tabular*}\vspace{-5pt}
}

\newcommand{\resumeSubSubheading}[2]{
    \begin{tabular*}{0.97\textwidth}{l@{\extracolsep{\fill}}r}
      \textit{\small#1} & \textit{\small #2} \\
    \end{tabular*}\vspace{-5pt}
}

\newcommand{\resumeSubItem}[2]{\resumeItem{#1}{#2}\vspace{-4pt}}

\renewcommand{\labelitemii}{$\circ$}

\newcommand{\resumeSubHeadingListStart}{\begin{itemize}[leftmargin=*]}
\newcommand{\resumeSubHeadingListEnd}{\end{itemize}}
\newcommand{\resumeItemListStart}{\begin{itemize}}
\newcommand{\resumeItemListEnd}{\end{itemize}\vspace{-5pt}}


%-------------------------------------------
%%%%%%  CV STARTS HERE  %%%%%%%%%%%%%%%%%%%%%%%%%%%%


\begin{document}

%----------HEADING-----------------
\begin{tabular*}{\textwidth}{l@{\extracolsep{\fill}}r}
  \textbf{Вадим Холодило} & Email: \href{mailto:vadimkholodilo@gmail.com}{vadimkholodilo@gmail.com}\\
  \href{https://www.linkedin.com/in/vadim-kholodilo-818a221a1/}{linkedin.com/in/vadim-kholodilo-818a221a1} \textbar\ Удалённо (готов к переезду) & Телефон: \href{tel:+79314112463}{+7-931-411-2463} \\
\end{tabular*}

%-----------SUMMARY-----------------
\section{Резюме}
\small{Backend-разработчик (C\#/.NET) с 4-летним опытом разработки микросервисов и REST API на .NET и ASP.NET Core. Работаю с MSSQL, PostgreSQL, Elasticsearch и RabbitMQ, применяю CQRS, пишу unit- и интеграционные тесты. Выстраивал CI/CD (GitLab CI, TeamCity, GitHub Actions), контейнеризацию на Docker, мониторинг и трассировку (Prometheus, Grafana, OpenTelemetry). Мигрировал legacy-код с .NET Framework на .NET, работал с AWS.}
\vspace{-8pt}

%-----------EDUCATION-----------------
\section{Образование}
\resumeSubHeadingListStart
\resumeSubheading
{Технологический университет Теннесси}{Куквилл, Теннесси}
{Бакалавр компьютерных наук, специализация «Кибербезопасность»}{Авг. 2019 --- Май 2024}
\resumeSubHeadingListEnd


%-----------EXPERIENCE-----------------
\section{Опыт работы}
\resumeSubHeadingListStart


\resumeSubheadingDesc
{Pyrus}{Москва, Россия}
{Backend-разработчик (C\#/.NET)}{Июн. 2024 --- Июл. 2026}
{Low-code платформа для автоматизации бизнес-процессов и рабочих процессов}
\resumeItemListStart
\resumeItem{Перенёс поиск и фильтрацию в модуле управления сотрудниками с frontend на Elasticsearch, обеспечив работу с более чем 40 000 записей о сотрудниках}
\resumeItem{Разрабатывал serverless-интеграции (боты) на AWS Lambda для клиентов из США, автоматизировав заполнение форм из Excel}
\resumeItem{Повысил стабильность автотестов в CI, сократив число нестабильных (flaky) тестов на 50\%}
\resumeItem{Исправил 1000+ багов в backend, повысив стабильность платформы}
\resumeItem{Провёл рефакторинг и миграцию Pyrus Sync (legacy-код) с .NET Framework 4.8 на .NET 10}
\resumeItem{Доработал мигратор для переноса организаций между RU и US кластерами: вместо копирования реализовал корректное преобразование сущностей (включая ID и историю тарифов) и исправление данных там, где это требовалось; обеспечил целостность данных за 10+ лет работы крупных организаций}
\resumeItem{Участвовал в code review; консультировал команду по Docker, контейнеризации и деплою}
\resumeItem{Доработал CI/CD пайплайны в GitLab CI и TeamCity для автоматизации релизов, сократив время сборки и тестирования на 30\%}
\resumeItem{Внедрил OpenTelemetry: трассировка и метрики упростили мониторинг и разбор проблем в продакшене}
\resumeItemListEnd

\resumeSubheadingDesc
{Build Logics}{Ноксвилл, Теннесси}
{Разработчик ПО (C\#/.NET)}{Нояб. 2023 --- Май 2024}
{CRM-системы для строительных и подрядных компаний}
\resumeItemListStart
\resumeItem{Интегрировал платёжную систему Stripe в CRM-систему на .NET 6 и ASP.NET Core, повысив безопасность и эффективность платежей}
\resumeItem{Настроил CI/CD пайплайны в GitHub Actions, оптимизировав доставку ПО и сократив время развёртывания на 50\%}
\resumeItem{Автоматизировал развёртывание окружений для разработки через Ansible и Vagrant (infrastructure as code), избавившись от ручной настройки окружений}
\resumeItem{Развернул Prometheus и Grafana для мониторинга серверов и приложений, настроил алерты}
\resumeItemListEnd

\resumeSubheadingDesc
{Alliance Bernstein}{Нашвилл, Теннесси}
{Разработчик ПО (C\#/.NET)}{Май 2022 --- Авг. 2023}
{Международная компания по управлению активами и финансовым услугам}
\resumeItemListStart
\resumeItem{Разработал и развернул 3 микросервиса с REST API для платформы управления заявками по облигациям на основе паттерна CQRS, повысив производительность системы на 25\%}
\resumeItem{Оптимизировал SQL-запросы в Microsoft SQL Server, увеличив скорость работы микросервиса Orders на 30\%}
\resumeItem{Настроил CI/CD пайплайны в Azure DevOps для новых и существующих микросервисов, сократив время развёртывания на 50\%}
\resumeItem{Реализовал событийно-ориентированные фоновые задачи на брокере сообщений RabbitMQ, разделив планирование и обработку для повышения масштабируемости и наблюдаемости}
\resumeItemListEnd

\begin{comment}
\resumeSubheading
{Технологический университет Теннесси}{Куквилл, Теннесси}
{Разработчик ПО (система индорной навигации и оповещения для людей с нарушениями зрения)}{Дек. 2020 -- Май. 2022}
\resumeItemListStart
\resumeItem{Разработал и внедрил надёжный backend и панель администрирования системы с использованием .NET, C\#, PostgreSQL, React, TypeScript, Material UI и Docker. Достиг рекордного срока вывода на рынок в 6 недель, опередив первоначальный план проекта на 30\%}
\resumeItem{Интегрировал систему с базой данных университета для получения данных о зданиях и аудиториях, полностью исключив ручной ввод данных и сэкономив 10 часов работы в неделю}
\resumeItem{Внедрил кэширование с использованием Redis, повысив производительность backend API на 50\%}
\resumeItem{Регулярно представлял результаты работы команде проекта, обеспечивая соблюдение сроков и выполнение всех требований}
\resumeItem{Успешно представил систему руководству университета, что привело к одобрению дальнейшего финансирования проекта в размере \$10 000}
\resumeItemListEnd
\resumeSubheading
{Технологический университет Теннесси}{Куквилл, Теннесси}
{Ассистент исследователя}{Авг. 2019 -- Дек. 2020}
\resumeItemListStart
\resumeItem{Разрабатывал исследовательские прототипы ПО на Python, JavaScript, C++ и Java}
\resumeItem{Измерял производительность и оценивал эффективность прототипов}
\resumeItem{Представлял результаты исследований на конференциях}
\resumeItem{Сотрудничал с другими исследователями в разработке инновационных решений}
\resumeItem{Стал соавтором и опубликовал 7 статей в авторитетных журналах, внося вклад в распространение значимых результатов исследований в области компьютерных наук}
\resumeItem{Наставлял и консультировал трёх новых ассистентов исследователя, способствуя их профессиональному росту в области компьютерных наук}
\resumeItemListEnd
\end{comment}
\resumeSubHeadingListEnd

%-----------PROJECTS-----------------
\begin{comment}
\section{Проекты}
\resumeSubHeadingListStart
\resumeSubItem{QuantSoftware Toolkit}
{Open source библиотека на Python для анализа финансовых данных и машинного обучения в финансовой сфере.}
\resumeSubItem{Github Visualization}
{Визуализация данных Git-лога с помощью D3 для анализа трендов проекта во времени.}
\resumeSubItem{Recommendation System}
{Рекомендательные системы для музыки и фильмов на основе коллаборативной фильтрации на публичных датасетах.}
\resumeSubItem{Mac Setup}
{Книга с пошаговыми инструкциями по настройке окружения разработчика на Mac OS.}
\resumeSubHeadingListEnd
\end{comment}

%
%--------PROGRAMMING SKILLS------------
\section{Ключевые навыки}
\resumeSubHeadingListStart
\item{
            \textbf{Языки и платформа}{: C\#, SQL, Python, .NET, .NET Core, ASP.NET Core, Entity Framework, REST API, микросервисы, CQRS, unit- и интеграционное тестирование, xUnit, MSTest, Testcontainers, Angular}
      }
\item{
            \textbf{Базы данных и очереди}{: MS SQL Server, PostgreSQL, Elasticsearch, RabbitMQ}
      }
\item{
            \textbf{DevOps и мониторинг}{: Docker, Linux, Git, CI/CD, GitLab CI, TeamCity, GitHub Actions, Azure DevOps, Ansible, Vagrant, AWS, Azure, OpenTelemetry, Prometheus, Grafana, Jira}
      }
\resumeSubHeadingListEnd


%-------------------------------------------
\end{document}
